\section{Courbure, torsion et action d’Einstein--Cartan--Holst}
\label{chap:gravedad-cartan-holst}

La structure de transport HMT distingue la variable de mémoire, l’holonomie et le défaut de fermeture. Sa réalisation gravitationnelle se formule au moyen d’une tétrade et d’une connexion. Cette section établit la géométrie du premier ordre et spécifie le type de l’application entre celle-ci et l’espace des états TPK\@.

\subsection{Identité discrète de Gauss–Stokes et transport de l’incidence vers la frontière}

Soit \(K\) un complexe cellulaire, soit \(W\) un groupe abélien et soit \(\Omega\in C^1(K;W)\) une cochaîne. Nous notons \(\delta:C^1(K;W)\to C^2(K;W)\) le cobord cellulaire et \(\langle\,\cdot\,,\,\cdot\,\rangle\) l’accouplement entre cochaînes à valeurs dans \(W\) et chaînes cellulaires entières.

\begin{theorem}[Identité cellulaire universelle de Stokes]
Pour toute \(2\)-chaîne finie \(z\in C_2(K;\mathbb Z)\),
\[
 \boxed{\langle\delta\Omega,z\rangle
 =\langle\Omega,\partial z\rangle.}
\]
\end{theorem}
\begin{proof}
C’est l’identité qui définit le cobord comme opérateur adjoint du bord cellulaire. La finitude de \(z\) garantit que les deux accouplements sont des sommes finies dans \(W\).
\end{proof}

\begin{corollary}[Gauss--Stokes sur une pseudovariété cellulaire]
Soit \(K\) une pseudovariété cellulaire régulière, bidimensionnelle, finie et compacte, orientée de manière cohérente, éventuellement à bord. Si \(z_K\) est la somme de ses \(2\)-cellules avec l’orientation cohérente, alors
\[
 \sum_{f\in K^{(2)}}(\delta\Omega)(f)
 =\sum_e[\partial z_K:e]\,\Omega(e).
\]
En particulier, la régularité permet d’identifier les incidences non annulées de \(\partial z_K\) aux arêtes du bord géométrique.
\end{corollary}
\begin{proof}
On applique le théorème à \(z_K\). L’orientation cohérente fait apparaître chaque arête intérieure deux fois dans \(\partial z_K\), avec des coefficients opposés ; les incidences non annulées forment la chaîne de bord.
\end{proof}

La version non abélienne ne s’obtient pas en remplaçant simplement les sommes par des produits. Dans cet ouvrage, elle n’est utilisée que lorsque les transports vers une base commune, un ordre des holonomies et la conjugaison correspondante ont été fixés. En dehors de ce sous-domaine, aucune identité de Stokes non abélienne n’est affirmée.

\subsection{Capacité aréale de la cellule qutrit : \(81\log 3\) et transport collectif d’incidence}

\begin{theorem}[Coupe minimale d’un cube discret]
Dans le graphe cubique \(N\times N\times N\) de capacité unitaire, toute coupe séparant deux faces opposées a une capacité au moins égale à \(N^2\), et il en existe une de capacité \(N^2\).
\end{theorem}
\begin{proof}
Il existe \(N^2\) chemins disjoints en arêtes qui traversent le cube dans la direction normale aux deux faces. Toute coupe doit les intercepter, donc sa capacité est au moins \(N^2\). Une couche transversale contient exactement \(N^2\) arêtes et atteint la borne.
\end{proof}

Pour \(N=9\), la coupe vaut 81 et
\[
 9^3=729=27^2.
\]
La dernière égalité permet de regrouper six trits sous forme d’un volume \(9^3\) ou d’une grille de cardinalité \(27^2\). C’est une bijection d’ensembles ; ce n’est pas un isomorphisme entre \((\ZZ/9\ZZ)^3\) et \((\ZZ/27\ZZ)^2\), dont les exposants diffèrent.

\subsection{Tétrade, connexion et \(2\) défauts}

Soient \(e^I\) une tétrade et \(\omega^{IJ}=-\omega^{JI}\) une connexion de Lorentz. Définissons
\[
 T_{\rm tor}^I=D_\omega e^I
 =\dd e^I+\omega^I{}_J\wedge e^J,
\]
\[
 F^{IJ}=\dd\omega^{IJ}+\omega^I{}_K\wedge\omega^{KJ},
 \qquad
 g=\eta_{IJ}e^I\otimes e^J.
\]
La décomposition
\[
 \omega^{IJ}=\mathring\omega^{IJ}(e)+K^{IJ}
\]
sépare la connexion de Levi--Civita et la contorsion. Alors
\[
 T_{\rm tor}^I=K^I{}_J\wedge e^J.
\]
Courbure et torsion sont des défauts différents : \(F\) mesure la fermeture du transport des vecteurs ; \(T_{\rm tor}\), la fermeture de la tétrade. La mémoire \(T\) du principe d’Ambivalence n’est pas non plus une torsion géométrique : seul un foncteur de réalisation peut les relier.

\subsection{Action du 1er ordre}

Soit \(M\) une variété lisse orientée de dimension quatre. Si \(\partial M\ne\varnothing\), nous supposons des variations à support compact dans l’intérieur ou, de manière équivalente pour le problème variationnel considéré, un terme de bord et des conditions aux limites annulant la contribution frontalière. Lorsque \(\Psi\) contient des champs spinoriels, nous supposons en outre une structure de spin et des champs compatibles avec elle. Soient
\[
 \Sigma^{IJ}=e^I\wedge e^J,
 \qquad
 (\star X)^{IJ}=\frac12\epsilon^{IJ}{}_{KL}X^{KL},
 \qquad
 \kappa=\frac{8\pi G}{c^4},
 \qquad
 \gamma\in\mathbb R\setminus\{0\}.
\]
L’action de Palatini--Cartan--Holst est
\[
\boxed{
 S[e,\omega,\Psi]
 =\frac1{2\kappa c}\int_M
 \Sigma_{IJ}\wedge\left(\star F^{IJ}+\frac1\gamma F^{IJ}\right)
 -\frac\Lambda{\kappa c}\int_M\operatorname{vol}_e
 +S_{\rm m}[e,\omega,\Psi].}
\]
La variation de cette action détermine les corrections torsionnelles \cite{Sciama1964,Kibble1961,Hehl1976,Holst1996}.

Définissons le courant de spin par
\[
 \delta_\omega S_{\rm m}
 =-\frac12\int_M\delta\omega^{IJ}\wedge\sigma_{IJ}.
\]

\begin{theorem}[Équation de Cartan--Holst]
Sous les hypothèses géométriques et de bord précédentes, la variation par rapport à \(\omega\) donne
\[
 D_\omega\!\left[(\star+\gamma^{-1}I)\Sigma^{IJ}\right]
 =\kappa c\,\sigma^{IJ},
\]
avec la normalisation globale fixée par la définition précédente de \(\sigma^{IJ}\).
\end{theorem}
\begin{proof}
On utilise \(\delta F^{IJ}=D_\omega\delta\omega^{IJ}\), on intègre par parties et l’on annule le coefficient de la variation arbitraire \(\delta\omega^{IJ}\). Avec la convention de signe fixée dans la définition de \(\sigma^{IJ}\), le terme de matière passe au membre de droite avec un signe positif. Le facteur \(1/(2\kappa c)\) produit le coefficient \(\kappa c\). Le terme cosmologique ne dépend pas de \(\omega\), et les conditions aux limites éliminent le terme engendré par l’intégration par parties.
\end{proof}

En signature lorentzienne, \(\star^2=-I\) sur les bivecteurs et, pour \(\gamma\in\RR\setminus\{0\}\),
\[
 (\star+\gamma^{-1}I)^{-1}
 =\frac{\gamma^2}{1+\gamma^2}(-\star+\gamma^{-1}I).
\]

\begin{corollary}[Découplage torsionnel]
Pour une cotétrade non dégénérée et \(\gamma\in\mathbb R\setminus\{0\}\),
\[
 \sigma^{IJ}=0\Longrightarrow T_{\rm tor}^I=0
 \Longrightarrow\omega=\mathring\omega(e).
\]
\end{corollary}

Pour expliciter la normalisation, définissons les sources normalisées par l’action au moyen de
\[
 \delta_g S_{\rm m}
 =-\frac12\int_M\operatorname{vol}_e\,
 \mathcal T^{S}_{\mu\nu}\,\delta g^{\mu\nu}.
\]
Lorsque les coordonnées de la tétrade ont la dimension d’une longueur, le tenseur énergie--impulsion physique conventionnel est \(T^{\rm phys}_{\mu\nu}:=c\,\mathcal T^{S}_{\mu\nu}\). Notons de même \(\mathcal U^{S,\rm spin}\) la contribution normalisée par l’action et \(U^{\rm phys,\rm spin}:=c\,\mathcal U^{S,\rm spin}\) sa convention physique. En éliminant algébriquement la connexion, l’équation métrique prend alors les deux formes équivalentes
\[
 G_{\mu\nu}+\Lambda g_{\mu\nu}
 =\kappa c\,
 \bigl(\mathcal T_{\mu\nu}^{S,\rm LC}
       +\mathcal U_{\mu\nu}^{S,\rm spin}\bigr)
 =\kappa\,
 \bigl(T_{\mu\nu}^{\rm phys,LC}
       +U_{\mu\nu}^{\rm phys,spin}\bigr),
\]
où la contribution du spin est quadratique dans le courant correspondant et dépend du couplage de matière. Le coefficient de cette contribution ne peut être fixé sans spécifier l’action fermionique et les conventions de dualité.

\subsection{Identité de Bianchi avec torsion et bilan de flux dans la connexion de Cartan}

Pour toute connexion,
\[
 D_\omega T_{\rm tor}^I=F^I{}_J\wedge e^J,
 \qquad
 D_\omega F^{IJ}=0.
\]
Après élimination de la torsion,
\[
\nabla^\mu(T_{\mu\nu}^{\rm phys,LC}
          +U_{\mu\nu}^{\rm phys,spin})=0.
\]
Ces identités sont des tests de cohérence obligatoires : tout terme gravitationnel HMT doit provenir d’une action compatible ou satisfaire le bilan de Noether correspondant.

\subsection{Application de réalisation et spécialisation du paramètre de Holst}

L’insertion définie dans \cref{mdg:chap:barbero} prend
\[
 \gamma=\Gamma_{6\to8}.
\]
Une fois cette substitution effectuée, toutes les conséquences variationnelles précédentes sont des théorèmes de l’action spécialisée. La sélection physique de cette action par le transport HMT s’exprime au moyen d’une application de réalisation
\[
 \mathfrak R_{\rm grav}:\mathcal X_{\rm TPK}
 \dashrightarrow\{(e,\omega,\Psi)\}
\]
dotée de trois propriétés : entrelacer le transport TPK avec celui de \(\omega\), identifier les classes d’holonomie et fournir les échelles de \(G\) et de \(\Lambda\). Ces propriétés fixent le domaine précis d’une réalisation gravitationnelle du formalisme.

Pour une cotétrade non dégénérée, \(\gamma\) réel non nul et une matière sans courant de spin, le secteur torsionnel se réduit à la formulation du premier ordre de la relativité générale avec la même constante cosmologique, sous les conditions aux limites déjà déclarées. Dans un fluide de spin, la loi d’échelle \(s^2\propto a^{-6}\) peut produire un rebond d’Einstein--Cartan ; obtenir sa densité et ses conditions initiales à partir de l’état TPK est une question distincte de la cohérence variationnelle de l’action.
